\section{P4 programmable data plane}
% TODO: put a section on bare-metal switches hardware %
Before turning to what P4 is and how the language itself is structured, it is worth establishing why a new data-plane programming model was needed at all. OpenFlow was built to reuse capability that switch vendors already had on-chip: McKeown et al.\ note that ``we exploit the fact that most modern Ethernet switches and routers contain flow-tables (typically built from TCAMs)'' used internally to implement functions such as firewalling, NAT, QoS, and statistics collection, and that OpenFlow simply exposes ``this common set of functions'' through a standard interface~\cite{mckeown2008openflow}. This was a deliberate compromise: rather than requiring new switch hardware, the design goal was a datapath ``amenable to high-performance and low-cost implementations,'' which meant ``forgoing the ability to specify arbitrary handling of each packet and seeking a more limited, but still useful, range of actions''~\cite{mckeown2008openflow}. In its original 2008 form this amounted to a single flow table matched against about a dozen header fields such as MAC address, IP address, protocol, and TCP/UDP ports~\cite{bosshart2014p4}.

% TODO: explain, in your own words, why this compromise became OpenFlow's fundamental limitation as the specification and hardware evolved.

\subsection{The fundamental limitation of OpenFlow}
That original dozen-field flow table did not remain fixed for long. The formal OpenFlow 1.0 specification followed the 2008 proposal in December 2009, and each successive revision of the OpenFlow 1.x line extended the matchable field set to accommodate new headers, so that the number of supported fields rose from 12 to 41 by 2013~\cite{bosshart2014p4}. Yet as data-center networking continued to evolve, operators kept demanding support for still newer encapsulation headers (e.g., VXLAN, NVGRE); rather than continuing to extend OpenFlow 1.x indefinitely, researchers proposed a strawman ``OpenFlow 2.0'' interface instead~\cite{bosshart2014p4}.

The push toward an OpenFlow 2.0 interface reflected two distinct pressures on any successor design: the need to keep pace with the rapid rate of change in networking requirements, and the need for a design that remained viable over the long term. The immediate problem was one of deployment speed. At the time, bare-metal switches still relied on a fixed-function pipeline hardcoded into the switch firmware to match packets, so supporting a new protocol required two slow steps in sequence. First, the protocol itself had to be agreed upon. Second, once agreed upon, the switch firmware itself had to be updated, since the firmware had no built-in intelligence to parse a raw packet format it had never been told about, and changing the firmware took still more time. Consequently, simply adding a new protocol to the specification was not enough on its own: it could still take vendors years to release a new switch model incorporating it, and for that hardware to then be adopted into the already-deployed network. The implication was that future switches should be able to define custom headers themselves, and that controller applications should likewise be able to understand these custom headers.

\subsection{P4 forwarding pipeline}
At the architecture level, researchers introduced PISA (Protocol Independent Switching Architecture) as the pipeline design that gives P4 its reconfigurability~\cite{systemsapproach}. Instead of a single fixed flow table, a PISA pipeline is built from three components---a parser, a multi-stage match-action pipeline, and a deparser---each configured by the P4 program itself rather than hardcoded into switch firmware. This generic, configurable structure is what lets the same abstract model describe packet processing across otherwise very different forwarding devices---Ethernet switches, load-balancers, and routers---and across different underlying technologies, whether a fixed-function switch ASIC, an NPU, a reconfigurable switch, a software switch, or an FPGA~\cite{bosshart2014p4}.


